\textbf{Partie A: étude du sens de variation d'une fonction}

\medskip

On considère la fonction $f$ définie sur $\R$ par 

\[f(x) = \dfrac{2x}{\sqrt{1 + x^2}}\]

\begin{enumerate}
\item Résoudre l'équation $f(x) = x$.
\item
	\begin{enumerate}
		\item On admet que la fonction $f$ est dérivable sur $\R$.
		
Vérifier que, pour tout réel $x$,on a $f'(x) = \dfrac{2}{(1+x^2)\sqrt{1 + x^2}}$
		\item En déduire le sens de variation de la fonction $f$ sur $\R$.
	\end{enumerate}
\end{enumerate}

\medskip

\textbf{Partie B : étude de la convergence d'une suite récurrente}

\medskip

La suite $(u_n)$ est définie par $u_0 = 1$ et, pour tout entier naturel $n,\: u_{n+1} = f(u_n)$.

\medskip

\begin{enumerate}
\item Démontrer par récurrence que, pour tout entier naturel $n$, on a $1 \leqslant u_n \leqslant u_{n+1} < \sqrt 3$.
\item En déduire que la suite $(u_n)$ converge et déterminer sa limite.
\item Le but de cette question est de retrouver par une autre méthode les résultats de la
question \textbf{2.} de la \textbf{partie B}.

Pour tout entier naturel $n$, on pose :

\[v_n = \dfrac{u_n^2}{3 - u_n^2}\]

On admet que la suite $(v_n)$ est bien définie.
	\begin{enumerate}
		\item Démontrer que la suite $(v_n)$ est une suite géométrique de raison 4 dont on précisera le premier terme.
		\item En déduire une expression de $v_n$ en fonction de $n$ puis que $u_n = \sqrt{\dfrac{1,5 \times 4^n}{1 + 0,5 \times 4^n}}$ pour tout entier naturel $n$.
		\item En déduire la limite de la suite $(u_n)$.
	\end{enumerate}
\end{enumerate}

\medskip

\textbf{Partie C : étude de la convergence de la somme de termes}

\medskip

Pour tout entier naturel $n$ non nul, on pose $S_n = u_0^2 + u_1^2 + \ldots + u_{n-1}^2$.

\medskip

\begin{enumerate}
\item Recopier et compléter le script Python ci-dessous afin que celui-ci permette de lister
les $p$ premiers termes de la suite $(S_n)$.

\begin{center}
\begin{ttfamily}
\begin{tabular}{|l|}\hline
\textbf{from} math \textbf{import}*\\
~\\
\textbf{def} termes(p) :\\
\qquad  u = \ldots\\
\qquad  S=0\\
\qquad  L= [ ]\\
\qquad  \textbf{for} i \textbf{in} range(p)\\
\qquad \qquad S=\ldots\\
\qquad \qquad u= \ldots\\
\qquad \qquad L.append(S) \\
\qquad \textbf{return} L\\ \hline
\end{tabular}
\end{ttfamily}
\end{center}

Remarque : on rappelle qu'en langage Python,

\begin{itemize}[label=$\bullet~$]
\item la commande L= [ ] crée une liste vide ;
\item la commande L.append (S) ajoute, à la fin de la liste L, l'élément supplémentaire S.
\end{itemize}

\item On rappelle que, pour tout entier naturel $k$, on a $1 \leqslant u_k \leqslant \sqrt 3$.

Démontrer que, pour tout entier naturel $n$ non nul, on a $n \leqslant S_n \leqslant 3n$.
\item En déduire les limites respectives de $S_n$ et de $\dfrac{S_n}{n^2}$ lorsque $n$ tend vers $+ \infty$.
\end{enumerate}

\medskip
